\chapter{Introdução}
\label{cap:introducao}

Os jogos digitais representam uma área de significativa importância na sociedade
contemporânea, com aplicações que abrangem desde o entretenimento até fins
educacionais. A indústria de jogos eletrônicos tem demonstrado um crescimento
notável, consolidando-se como um mercado de alcance global sem precedentes
\cite{wijman2023newzoo}.

Nesse contexto, diversos jogos do gênero \textit{roguelike} têm alcançado grande
sucesso, com títulos notáveis como \textit{Hollow
    Knight}\footnote{Página da Steam do Hollow Knight (Acesso em: 23 jan. 2026): https://store.steampowered.com/app/367520/Hollow\_Knight/},
\textit{Undermine}\footnote{Página da Steam do UnderMine (Acesso em: 23 jan. 2026): \url{https://store.steampowered.com/app/656350/UnderMine/}}
e \textit{Dead
    Cells}\footnote{Página da Steam do Dead Cells (Acesso em: 23 jan. 2026): \url{https://store.steampowered.com/app/588650/Dead\_Cells/}}
destacando-se pela capacidade de proporcionar experiências desafiadoras e
imersivas. O gênero \textit{roguelike}, um subgênero de \gls{RPG},
caracteriza-se fundamentalmente pela exploração de masmorras geradas
aleatoriamente, combate tipicamente baseado em turnos e a morte permanente do
personagem, que impõe o recomeço do jogo após uma derrota
\cite{laurindo2024relatorio}.

Um elemento fundamental que define a experiência em muitos jogos
\textit{roguelike} é a \gls{PCG}. \gls{PCG} é uma técnica que utiliza algoritmos
para criar automaticamente conteúdos de jogo, como mapas e cenários, sendo
fundamental para garantir que cada sessão de jogo seja única e ofereça alta
rejogabilidade \cite{sepanmaa2024procedural}. Embora a \gls{PCG} seja poderosa para
gerar uma vasta quantidade de conteúdo, um de seus desafios históricos reside na
dificuldade de controlar o resultado. Tradicionalmente, os desenvolvedores
manipulam parâmetros complexos e abstratos para guiar o processo gerador, o que
muitas vezes torna a criação de conteúdo com intenções específicas uma tarefa
pouco intuitiva e iterativa \cite{lazaridis2024creating}.

A busca por maior controle na \gls{PCG} tem sido um campo ativo de
pesquisa.  Contudo, gerar conteúdo que reflita intenções específicas e
restrições de alto nível permanece um desafio. Recentemente, os avanços em
\gls{LLM}s abriram novas e promissoras formas para abordar este problema. A
capacidade inerente dos \gls{LLM}s de compreender e processar linguagem natural
oferece uma interface fundamentalmente nova para a geração de conteúdo: em vez
de ajustar parâmetros, os desenvolvedores ou mesmo os jogadores poderiam
simplesmente descrever o conteúdo desejado textualmente.

Essa abordagem, frequentemente denominada \textit{"text-to-level"} ou
\textit{"text-to-map"}, transforma o \gls{LLM} em um tradutor de intenções
humanas para estruturas de jogo concretas. Um trabalho pioneiro nesta área,
\textit{MarioGPT}, demonstrou a viabilidade de treinar um \gls{LLM} para gerar
níveis de \textit{Super Mario Bros} a partir de descrições simples como "muitos
canos, poucos inimigos, elevação baixa"\ \cite{sudhakaran2023mariogpt}. Este
estudo provou que os \gls{LLM}s podem não apenas gerar sequências textuais, mas
também produzir conteúdo estruturado e funcional para jogos, abordando um dos
principais desafios da \gls{PCG}.

Inspirado por esses avanços, este trabalho investiga a aplicação dessa abordagem
no domínio dos jogos \textit{roguelike}. Esse gênero é caracterizado por mapas
com restrições estruturais, como a conectividade entre salas, a disposição de
corredores e a distribuição lógica de inimigos e tesouros. Sendo assim, o
objetivo principal é investigar como os LLMs podem traduzir uma descrição
textual em um mapa funcional, garantindo que a estrutura resultante se adéque a
essas convenções.

Este estudo busca desenvolver um sistema que utilize um \gls{LLM} para
interpretar descrições textuais e produzir uma representação de mapa
estruturada, que possa ser posteriormente renderizada por um motor de jogo. A
pesquisa se concentrará na eficácia de técnicas de engenharia de \textit{prompt}
para guiar o modelo, garantindo que os mapas gerados não apenas correspondam à
descrição, mas também sejam funcionalmente válidos, respeitando as restrições
estruturais do gênero mencionadas anteriormente (como conectividade de salas e
disposição de corredores).  Acredita-se que tal abordagem oferece um método de
criação de conteúdo procedural mais intuitivo e expressivo em comparação com
técnicas algorítmicas tradicionais, como Perlin Noise \cite{perlin1985image} e
Random Walk \cite{pearson1905problem}, contribuindo com conhecimento técnico e
metodológico.

\section{Contribuições}
\label{sec:contribuicoes}

Considerando os resultados obtidos e a validação do sistema proposto, este
trabalho apresenta contribuições técnicas e práticas para a área de \gls{PCG}
assistida por Inteligência Artificial. As principais contribuições podem ser
sintetizadas em três pilares:

\begin{itemize}
    \item \textbf{Arquitetura Híbrida de Geração "Text-to-Map":} O
          desenvolvimento de uma metodologia que desacopla a responsabilidade criativa
          da responsabilidade estrutural. Esta abordagem mitiga as limitações de
          raciocínio espacial inerentes aos \gls{LLM}s, delegando a topologia a
          algoritmos determinísticos enquanto utiliza o modelo de linguagem para
          enriquecer semanticamente o jogo (narrativa, inimigos e itens). O sistema
          materializa essa criação através de arquivos estruturados em \gls{JSON},
          garantindo interoperabilidade com motores de jogo (um exemplo de um artefato
          JSON gerado encontra-se disponível no Apêndice~\ref{ap:exemplo_json}).

    \item \textbf{Dataset Curado para RAG em Jogos:} A implementação de um
          \textit{pipeline} técnico que integra \textit{Structured Outputs} e
          \gls{RAG}. Neste contexto, a curadoria e anotação manual de um
          \textit{dataset}, composto por 489 pares de descrições e imagens, constitui
          uma contribuição técnica relevante deste trabalho, estabelecendo uma base
          sólida para a associação entre linguagem natural e ativos de \textit{pixel
              art}.

    \item \textbf{Ferramentas de Código Aberto:} A disponibilização de dois
          artefatos de \textit{software} funcionais e integrados, que servem como
          prova de conceito e base para pesquisas futuras. Os códigos-fonte do projeto
          foram tornados públicos nos seguintes repositórios:
          \begin{itemize}
              \item \textbf{Gerador de Assets (Python):} \url{https://github.com/GusGurgel/tcc-roguelike-assets-generator}
              \item \textbf{Gerador de Mapas (Godot):} \url{https://github.com/GusGurgel/tcc-roguelike}
          \end{itemize}
\end{itemize}